\section{Introduction}
The "Best Bike Paths" (BBP) system is a comprehensive software platform commissioned by a cyclist association to serve both its members and the broader cycling community. The system addresses the need for a collaborative repository of bicycle paths, facilitating their creation, maintenance, and utilization. It combines personal activity tracking with crowdsourced territorial mapping, leveraging both manual user reporting and automated data collection via mobile device sensors to identify path conditions and obstacles.
\subsection{Purpose}
The purpose of this document is to provide a comprehensive architectural and detailed design for the "Best Bike Paths" (BBP) system. It builds upon the requirements defined in the Requirement Analysis and Specification Document (RASD) and aims to guide the development team through the implementation, testing, and integration phases. Specifically, this document defines the system's high-level architecture, the decomposition into logical components, the user interface design, and the mapping between requirements and software modules.

From a business perspective, the purpose of the product is to enhance the cycling experience by enabling users to track their performance and fostering a safer environment through community-driven updates on path quality and obstacles.

\subsection{Scope}
The BBP system involves interactions between two main user categories: Registered Users and Unregistered Users (Explorers), interacting with the system primarily through a mobile application, supported by backend services.

\textbf{Registered Users} can access the mobile application to record trips, utilizing the "Automated Mode" to detect anomalies (e.g., potholes) via device sensors (accelerometer/gyroscope). They can verify detected data, manually report path status, and manage their personal history and statistics. They are also responsible for publishing private data to the shared inventory.

\textbf{Unregistered Users} can access the "Explore" section to plan routes. They input an origin and destination, and the system calculates optimal paths based on a "Path Score," which aggregates distance and surface quality derived from community data.

\textbf{The System }calculates trip statistics, enriches data with meteorological information from external services, and manages the calculation of path scores based on valid reports.


\subsection{Definitions, Acronyms, Abbreviations}

\subsubsection{Definitions}
\begin{itemize}
    \item \textbf{Registered User:} A user with a platform account who can record trips, report anomalies, and access personal history.
    \item \textbf{Unregistered User:} A user who can only view maps and plan routes without saving data or tracking activities.
    \item \textbf{Trip:} A recorded cycling session containing geographical waypoints, duration, performance statistics, and weather data.
    \item \textbf{Automated Mode:} A feature that detects cycling activity and records anomalies via mobile sensors without manual intervention.
    \item \textbf{Manual Mode:} A data entry method where the user manually inputs status details or obstacles for a specific path segment.
    \item \textbf{Anomaly:} An irregularity on the road surface, such as a pothole, detected by the smartphone's accelerometer and gyroscope.
    \item \textbf{Confirmation:} The action of validating a detected anomaly, marking it as a verified obstacle or a false positive.
    \item \textbf{Path Score:} A numerical value used to rank paths based on maintenance status and travel efficiency.
    \item \textbf{Thick Client:} A mobile application that hosts significant business logic, such as real-time sensor processing and offline buffering.
    \item \textbf{Fat Server:} A backend server that handles long-term persistence, identity management, and intensive calculations like path scoring.
    \item \textbf{Tier:} A physical or logical layer in the system architecture where specific processing or data storage occurs.
    \item \textbf{RESTful:} An architectural style used for communication between the mobile client and backend server via web services.
\end{itemize}

\subsubsection{Acronyms}
\begin{itemize}
    \item \textbf{API:} Application Programming Interface.
    \item \textbf{BBP:} Best Bike Paths.
    \item \textbf{DBMS:} Database Management System.
    \item \textbf{DD:} Design Document.
    \item \textbf{GDPR:} General Data Protection Regulation.
    \item \textbf{GPS:} Global Positioning System.
    \item \textbf{HTTPS:} Hypertext Transfer Protocol Secure.
    \item \textbf{JSON:} JavaScript Object Notation, used for RESTful data exchange.
    \item \textbf{JWT:} JSON Web Token, used for secure session authentication.
    \item \textbf{MVC:} Model-View-Controller architectural pattern.
    \item \textbf{RASD:} Requirements Analysis and Specification Document.
    \item \textbf{UI:} User Interface.
\end{itemize}

\subsubsection{Abbreviations}
\begin{itemize}
    \item \textbf{G$n$:} Goal number $n$
    \item \textbf{R$n$:} Requirement number $n$
    \item \textbf{PR$n$:} Performance Requirement number $n$
    \item \textbf{CON$n$:} Constraint number $n$
    \item \textbf{DC$n$:} Design Constraint number $n$
    \item \textbf{SSA$n$:} Software System Attribute number $n$
\end{itemize}

\subsection{Revision History}
\begin{itemize}
    \item January 7, 2026: version 1.0, first release
\end{itemize}

\subsection{Reference Documents}
\begin{itemize}
    \item Specification document: "Assignment RDD AY 2025-2026"
    
    \item UML official specification: https://www.planttext.com/
    
    \item Sequence Diagram specification: https://www.planttext.com/
    
    \item User Interface Design Figma project:
    \href{https://www.figma.com/design/Eke18Kn9sdAbKgfQLGRc6u/BBP?node-id=0-1&p=f&t=du3CXIlhp8B1eogT-0}{FigmaProject-BBP}
    
    \item Deployment Diagram specification: 
    \href{https://online.visual-paradigm.com/diagrams/tutorials/deployment-diagram-tutorial/}{Visual-Paradigm-deployment-diagram}
    
\end{itemize}

\subsection{Document Structure}
\begin{itemize}
    \item \textbf{Section 1: Introduction}\\
    This section sets the context for the entire document. It outlines the scope and purpose of the design, while providing a comprehensive glossary of definitions, acronyms, and abbreviations to ensure a consistent understanding of the terminology used throughout the project.

    \item \textbf{Section 2: Architectural Design}\\
    Directed primarily at the development team, this section provides a deep dive into the system's technical framework. It details the chosen architectural styles and patterns, decomposing the system into its core components and layers. Additionally, it presents the runtime and deployment views to illustrate how software modules interact and map to the hardware infrastructure.

    \item \textbf{Section 3: User Interface Design}\\
    This section focuses on the user experience (UX) and visual aspects of the application. It presents a collection of interface mockups and navigation flowcharts, serving as a guideline for the frontend design of the mobile application and ensuring alignment with user interaction requirements.

    \item \textbf{Section 4: Requirements Traceability}\\
    This section ensures consistency between the analysis and design phases. It provides a detailed mapping that links the functional requirements and goals defined in the RASD to the specific architectural modules and components described in this document, verifying that all requirements are addressed.

    \item \textbf{Section 5: Implementation, Integration and Test Plan}\\
    The final section outlines the strategy for bringing the system to life. It describes the planned sequence for implementation and integration of the various modules. Furthermore, it defines the testing protocols necessary to verify that the core functionalities perform as expected within the integrated environment.
\end{itemize}